% TikZ 图：安培力的微观来源
% Day 22：洛伦兹力
\begin{tikzpicture}[scale=0.9,>=Latex,line cap=round]
  \path[use as bounding box] (-0.35,-2.25) rectangle (7.85,2.30);
  \tikzset{
    wire/.style={thick,fill=gray!12},
    carrier/.style={circle,fill=red!12,draw=red!80!black,inner sep=1.2pt},
    vec/.style={->,very thick}
  }

  \draw[wire] (0,0) rectangle (7.0,0.90);
  \draw[<->,thick] (0,1.05) -- (7.0,1.05);
  \node[above] at (3.50,1.05) {$L$};
  \node[font=\small] at (3.50,-0.28) {正载流子等效模型};

  \foreach \x in {0.65,1.45,...,6.65} {
    \foreach \y in {-0.75,-1.35} {
      \draw[blue!70!black,thick] (\x-0.09,\y-0.09) -- (\x+0.09,\y+0.09);
      \draw[blue!70!black,thick] (\x-0.09,\y+0.09) -- (\x+0.09,\y-0.09);
    }
  }
  \node[blue!70!black,font=\small] at (7.35,-1.05) {$B$ 向里};

  \foreach \x/\y in {0.75/0.28,1.75/0.62,2.75/0.28,3.75/0.62,4.75/0.28,5.75/0.62} {
    \node[carrier,font=\scriptsize] at (\x,\y) {$+$};
    \draw[vec,red] (\x+0.13,\y) -- ++(0.44,0);
  }
  \node[red,font=\small] at (3.45,0.36) {$v_d$};

  \draw[vec,red] (0.45,1.55) -- (2.10,1.55);
  \node[red,above] at (1.28,1.55) {$I=n|q|Sv_d$};

  \draw[vec,orange!85!black] (6.35,1.18) -- (6.35,1.82);
  \node[orange!85!black,right,font=\small] at (6.35,1.52)
    {$F_{\text{安}}=BIL$};

  \node[font=\small,align=center] at (3.50,-1.95)
    {载流子所受磁场力传给晶格，宏观上表现为导线所受安培力};
\end{tikzpicture}
